\chapter{Le articolazioni dell'arto inferiore}

% ---------------------------------------------------------------------------
\section{Generalità e articolazioni del bacino}

Lo scheletro dell'arto inferiore sostiene e trasmette a terra il peso, permette la deambulazione
bipede e l'equilibrio.
\begin{itemize}
  \item \textbf{cingolo pelvico}: raccordo con lo scheletro assile (ischio, pube, ileo);
  \item \textbf{coscia}: femore; \textbf{gamba}: tibia + fibula; \textbf{piede}: tarso (7 ossa
    brevi), metatarso (5 lunghe), falangi (3 per dito, 2 per l'alluce);
  \item articolazioni: bacino, coxo-femorale, ginocchio, tibio-fibulare, piede.
\end{itemize}

\textbf{Bacino}: sostegno strutturale, vincolo dinamico tronco--arto, ancoraggio muscolare,
contenitore di organi (digerente, urinario, genitale). Articolazioni:
\begin{itemize}
  \item \textbf{sacroiliaca}: tra osso dell'anca e sacro;
  \item \textbf{sinfisi pubica}: tra i due rami pubici.
\end{itemize}
Legamenti della pelvi: anteriori (longitudinale anteriore, ileolombare, sacroiliaci anteriori,
inguinale, sacrotuberoso, sacrospinoso, membrana otturatoria) e posteriori (sacroiliaci interossei e
posteriori, tesi verso i processi spinosi lombari, in particolare L4).

% ---------------------------------------------------------------------------
\section{Articolazione coxo-femorale}

Tra acetabolo e testa del femore $\rightarrow$ \textbf{enartrosi} (a sfera, come la glenomerale), 3
assi, 6 movimenti.
\begin{itemize}
  \item cinque legamenti: ileofemorale, ischiofemorale, pubofemorale, rotondo del femore, acetabolare
    trasverso;
  \item \textbf{legamento trasverso dell'acetabolo}: chiude l'incisura acetabolare (margine inferiore
    della fossa);
  \item \textbf{capsula articolare}: sul contorno/labbro dell'acetabolo e sul femore (linea
    intertrocanterica avanti, collo anatomico dietro);
  \item \textbf{ileo-femorale}: dal bordo anteriore del coxale alla linea intertrocanterica; si
    divide in fascio pretrocanterico (il più potente dell'anca) e fascio al piccolo trocantere;
  \item \textbf{legamento della testa femorale} (rotondo): benderella robusta ($\sim$3--3,5 cm), dalla
    fossetta della testa alla fossa dell'acetabolo; contiene la branca posteriore dell'arteria
    otturatoria (ma la testa è vascolarizzata anche dalle arterie capsulari: circonflessa mediale/
    laterale, profonda del femore); intrarticolare, avvolto dalla sinovia.
\end{itemize}
\rubrica{Movimenti}
\begin{itemize}
  \item \textbf{flessione}: $\sim$90° a ginocchio esteso, 120° a ginocchio flesso, fino a 145° in
    massima mobilità;
  \item \textbf{estensione} $\sim$20°; \textbf{abduzione} $\sim$30°;
  \item combinati: estensione-adduzione, flessione-adduzione.
\end{itemize}

% ---------------------------------------------------------------------------
\section{Articolazione del ginocchio}

Complesso articolare, \textbf{ginglimo angolare} (a troclea, 1 asse, 2 movimenti). Coinvolge condili
femorali, condili tibiali (cavità glenoidee della tibia) e faccia posteriore della patella.

\subsection{Menischi e lesioni meniscali}
\textbf{Menischi} (mediale e laterale): aumentano la congruenza e distribuiscono le forze;
accompagnano i condili femorali sulle cavità tibiali.
\begin{itemize}
  \item \textbf{legamento trasverso del ginocchio}: collega le corna anteriori dei due menischi (poi
    alla rotula); il menisco mediale è connesso anche dal legamento meniscofemorale posteriore;
  \item durante estensioni rapide (calcio) i menischi possono essere pizzicati $\rightarrow$
    \textbf{lesioni meniscali}, soprattutto il \textbf{mediale} (meno mobile);
  \item clinica: dolore in rotazione laterale = menisco laterale; in rotazione mediale = menisco
    mediale;
  \item tipi: radiale, a flap, longitudinale, degenerativa; quadri: manico di secchio, frattura
    radiale del corno posteriore.
\end{itemize}

\subsection{Mezzi di unione e borse}
\begin{itemize}
  \item \textbf{capsula fibrosa}: riveste l'articolazione;
  \item \textbf{membrana sinoviale}: internamente; si interrompe ai menischi; posteriormente circonda
    i crociati $\rightarrow$ intracapsulari ma fuori dalla cavità; i collaterali sono extracapsulari;
  \item borse: \textbf{sovrapatellare} (femore--quadricipite), \textbf{prepatellare} (cute--patella),
    \textbf{infrapatellare} (tibia--legamento patellare).
\end{itemize}

\subsection{Stabilità e legamenti}
Articolazione meccanicamente debole $\rightarrow$ stabilizzatori \textbf{attivi} (muscoli/tendini) e
\textbf{passivi} (legamenti). Gruppi: anteriori, laterali, posteriori, intra-articolari.
\begin{itemize}
  \item \textbf{patellare} (anteriore): dal contorno inferiore della rotula alla tuberosità tibiale
    (tratto terminale del tendine quadricipitale); la femororotulea è un'\textbf{artrodia};
  \item \textbf{retinacoli} mediale e laterale della patella: dalle aponeurosi dei vasti $\rightarrow$
    stabilizzano la patella;
  \item \textbf{collaterale mediale/tibiale}: epicondilo mediale $\rightarrow$ faccia mediale della
    tibia; impedisce la lateralità esterna (valgo);
  \item \textbf{collaterale laterale/fibulare}: epicondilo laterale $\rightarrow$ testa della fibula;
    impedisce la lateralità interna (varo); entrambi tesi a ginocchio esteso;
  \item \textbf{popliteo obliquo} (espansione del semimembranoso) e \textbf{popliteo arcuato}
    (posteriori);
  \item \textbf{crociato anteriore (LCA)}: area intercondiloidea anteriore della tibia $\rightarrow$
    condilo laterale del femore; \textbf{crociato posteriore} (più robusto): area intercondiloidea
    posteriore $\rightarrow$ condilo mediale; sempre tesi, impediscono i \textbf{movimenti a cassetto}
    (scorrimento antero-posteriore della tibia).
\end{itemize}
Intrarotazione della gamba flessa limitata dai crociati; extrarotazione arrestata dai collaterali.

\subsection{Nota clinica: rottura dei crociati}
\begin{itemize}
  \item rottura $\rightarrow$ destabilizzazione (fenomeno del cassetto anteriore/posteriore);
  \item LCA: $\sim$10 volte più frequente del posteriore; meccanismo tipico = intrarotazione forzata
    con gamba bloccata; spesso da colpo diretto (calcio);
  \item colpo laterale a ginocchio esteso con piede bloccato $\rightarrow$ \textbf{``infelice
    triade''} (LCA + collaterale mediale + menisco mediale);
  \item crociato posteriore: più raro, da traumi sportivi o \textbf{trauma da cruscotto} (impatto del
    ginocchio nell'incidente stradale).
\end{itemize}

\subsection{Cinematica}
Flessione ed estensione della gamba sulla coscia.

% ---------------------------------------------------------------------------
\section{Articolazione tibio-fibulare}

\textbf{Membrana interossea} della gamba: tra i margini interossei di tibia e perone $\rightarrow$
stabilizza le tibio-fibulari e separa i muscoli anteriori dai posteriori. Distalmente, malleolo
mediale e laterale formano la \textbf{pinza malleolare}.
\begin{itemize}
  \item \textbf{prossimale}: rinforzata dai legamenti della testa della fibula (anteriori/posteriori);
  \item \textbf{distale}: \textbf{sinartrosi} tra le estremità distali (incisura fibulare della tibia
    con superficie rugosa/piana della fibula).
\end{itemize}

% ---------------------------------------------------------------------------
\section{Le articolazioni del piede}

Comprendono caviglia, tarso, metatarso, dita.
\begin{itemize}
  \item \textbf{tibio-tarsica} (talo-crurale): \textbf{ginglimo angolare} (1 asse) $\rightarrow$
    flessione dorsale e plantare; formata da pinza malleolare + astragalo;
  \item \textbf{intertarsali}: talo-calcaneo-navicolare (anteriore) e talo-calcaneale (posteriore,
    trocoide); \textbf{calcaneo-cuboidea} = a sella;
  \item \textbf{tarsometatarsali}: artrodie (fila distale del tarso -- basi dei metatarsali); piccoli
    movimenti + modifiche della volta plantare;
  \item \textbf{intermetatarsali}: tra le basi dei metatarsali (le prime due unite da legamento
    interosseo); scivolamento;
  \item \textbf{metatarsofalangee}: a condilo; flesso-estensione delle dita;
  \item \textbf{interfalangee}: a troclea; due per dito (una per l'alluce); flesso-estensione.
\end{itemize}
